\section{Data details}\label{app:data}
% =============================================================================================
% Owner: data writer. Tables tab:data-universe, tab:data-anomalies and tab:data-panel are
% written by paper/scripts/data_appendix.py (synthval.data on the raw file;
% results/e3/meta_fit.json). tab:e3-fit-diagnostics-full and tab:e3-heston-moments are generated
% by experiments/e3_sp500.py (part fit). Other numbers: results/e3/e3_fit_*.csv,
% results/e4/universes.csv. Fit diagnostics of the E4 generators: app_mk.tex.
% =============================================================================================

This appendix documents the construction of the ticker universe (\cref{sec:data-panel}) and
gives further diagnostics of the generator fits (\cref{sec:data-generators}). The fitting
algorithms themselves are described in \cref{app:impl}.

\paragraph{Universe.}
The Bloomberg export contains \texttt{PX\_LAST} (closing price) and \texttt{PX\_VOLUME} for the
503 tickers that were S\&P~500 constituents when the file was extracted; its requested end date
is 2025-06-17 and its last trading day 2025-06-16. Only prices enter the analysis; volume
serves the anomaly screen.
\Cref{tab:data-universe} gives the steps from the file to the panel. The availability rule
(a price on the first and on the last trading day, internal gaps of at most three days
forward-filled, nothing back-filled) removes the 22 tickers whose history starts after
2018-01-02. None of the remaining 481 tickers has a missing price, so the forward fill is never
used.

\begin{table}[tbp]
  \centering
  \caption{Construction of the ticker universe from the Bloomberg export.}
  \label{tab:data-universe}
  \small
  \fittable{manual_data_universe}
  \tablenote{Computed from the raw file by \texttt{synthval.data} (table written by
    \texttt{paper/scripts/data\_appendix.py}). Complete price history: a price on the first and
    the last trading day and no internal gap longer than three days. Anomaly screen: see
    \cref{tab:data-anomalies}.}
\end{table}

\paragraph{Anomaly screen.}
Two complete-history tickers are data errors. The screen that finds them is a rule, not a
hand-picked list: a ticker is flagged if one daily log-return exceeds 1 in absolute value, which
would mean a price ratio beyond $e\approx2.72$ in one day and points to an unadjusted corporate
action, or if it has at least 20 days with a price but zero volume, which points to stale quotes.
On the current file the rule flags exactly KDP and AMCR (\cref{tab:data-anomalies}). The
\texttt{PX\_LAST} series of KDP is not adjusted for the special dividend paid when the Keurig Dr
Pepper merger closed: the price falls from 123.66 on 2018-07-09 to 22.19 on 2018-07-10, on a
volume of 23.6 million shares, which leaves one spurious log-return of $-1.718$. Before its NYSE
listing on 2019-06-11, AMCR is a stale quote: 361 days with only 71 distinct prices and zero
volume on 232 of them; 298 of its 1258 \TRAIN{} returns are exactly zero (268 of them before
the listing), with runs of up to 18 days. Both thresholds are far from the values of the
retained tickers: the largest absolute log-return among them is 0.774 (APA on 2020-03-09, the
day of the oil price collapse), and none has more than one zero-volume day. One retained ticker
is unusual without being an error: SMCI has 117 zero \TRAIN{} returns, because its
split-adjusted \TRAIN{} prices (between 1.16 and 9.39) are quoted to two decimals; it is not in
the main panel but appears in some E4 universes.

\begin{table}[tbp]
  \centering
  \caption{Anomaly screen: the two excluded tickers and, for reference, the extreme and median
    values over the 479 retained tickers.}
  \label{tab:data-anomalies}
  \small
  \fittable{manual_data_anomalies}
  \tablenote{Largest absolute daily log-return over 2018-01-03 to 2025-06-16 and its date;
    days with a price but zero volume over the whole file; exactly zero returns and their
    longest run in \TRAIN{} (1258 returns). A ticker is flagged if the largest absolute
    log-return exceeds 1 or if it has at least 20 zero-volume days. Computed from the raw file
    by \texttt{synthval.data} (table written by \texttt{paper/scripts/data\_appendix.py}).}
\end{table}

\paragraph{Main panel and E4 universes.}
The main panel (\cref{tab:data-panel}) is drawn once: the 481 complete-history tickers, sorted
alphabetically, are permuted with a generator seeded by 20251114, and the first 100 tickers of
the permutation that are not excluded are kept. The draw without the exclusion would have
contained KDP and AMCR; with it they are replaced by the next two tickers of the permutation,
PCAR and TMUS, and the other 98 tickers are unchanged. The E4 universes are 30 independent
draws of 30 names without replacement from the 479 tickers of the default universe, each with
its own seed (\cref{sec:mk-design}); universes may overlap, and together they contain 400
distinct names. Fit diagnostics of the E4 generators are in \cref{app:mk-generators}.

\begin{table}[tbp]
  \centering
  \caption{The 100 tickers of the main panel (E1, E3, E5).}
  \label{tab:data-panel}
  \small
  \fittable{manual_data_panel}
  \tablenote{Seeded draw from the 479-ticker default universe (seed 20251114); read from the
    metadata of the E3 fits (\texttt{results/e3/meta\_fit.json}) and checked against
    \texttt{synthval.data.panel\_tickers}.}
\end{table}

\paragraph{Univariate fits.}
\Cref{tab:e3-fit-diagnostics-full} is the full version of \cref{tab:e3-fit-diagnostics}, with
the drifts, the \GARCHt{} intercept, optimiser success and fit times. The parameters at an
optimiser bound belong to the following tickers. The Merton intensity sits at its upper bound
of 0.5 jumps per day for GOOG, MTB and NVDA, whose fits describe frequent small jumps rather than
rare large ones. The only \GARCHt{} fit at the persistence cap is CTRA ($a=0.033$, $b=0.962$);
with variance targeting, the unconditional variance of every \GARCHt{} fit equals the \TRAIN{}
sample variance. The 33 Heston fits with a parameter at a bound have $\rho_v=-0.95$ (20
tickers), $c=\sqrt{1.5}$ (11) or $\kappa_v=0.5$ (5), the upper bound of the mean-reversion rate.
The stationary-bootstrap block length ranges from 15.8 to 60.4 days across tickers, and the
rule never falls back to its default.

\begin{table}[tbp]
  \centering
  \caption{Fits of the univariate generators on \TRAIN{}, full version (E3, 100 tickers, one fit
    per ticker and generator).}
  \label{tab:e3-fit-diagnostics-full}
  \small
  \fittable{e3_fit_diagnostics}
  \tablenote{Maximum likelihood (GBM, Merton, \GARCHt{}) and simulated method of moments
    (Heston) on the 1258 \TRAIN{} returns of each ticker, daily units. Median and interquartile
    range [25\%, 75\%] across the 100 tickers. Symbols: $\lambda=\lambda_J$; $\alpha,\beta$ are
    the GARCH coefficients $a,b$ of \cref{sec:data-generators} ($\omega$ is tied to the sample
    variance by variance targeting); $\kappa,\theta,\xi,\rho$ are
    $\kappa_v,\theta_v,\xi_v,\rho_v$. At bound: number of tickers whose estimate sits at an
    optimiser bound; the last row of each block counts tickers with at least one parameter at a
    bound. Bootstrap: mean block length in days.}
\end{table}

\paragraph{Heston simulation.}
Measure the variance in units of $\theta_v$, so that it has unit mean, and let $A=1-e^{-\kappa_v}$
and $B=e^{-\kappa_v}$. Given today's value $v$, tomorrow's variance has conditional mean
$m=A+vB$ and conditional variance $s^2=c^2(2vAB+A^2)$ (the transition moments of the square-root
process over one day, with $\xi_v^2=2\kappa_v\theta_vc^2$). Since $m^2=A^2+2vAB+v^2B^2$, the
squared coefficient of variation of the transition satisfies $\psi=s^2/m^2\le c^2$ for every $v$.
The QE scheme of \citet{andersen_2008_simple} draws tomorrow's variance as a moment-matched
multiple of a squared shifted Gaussian variable when $\psi\le\psi_c$ and from a mixture of a point
mass at zero and an exponential law otherwise; we use $\psi_c=1.5$. With $c\le\sqrt{1.5}$ the
first branch is used on every day, so the simulated variance is positive almost surely; no day
of the calibration paths of any of the 100 fits has zero variance. For larger $c$ the second
branch puts mass at zero, and a day that starts and ends at zero variance has a constant return.

\Cref{tab:e3-heston-moments} compares the data and simulated values of the seven moments of the
Heston calibration and of the kurtosis, which is not targeted. The variance, the tail quantile
ratio, the autocorrelation of $r_t^2$ at lag 20 and the leverage correlation are matched closely
(median absolute relative error at most 0.21). The autocorrelations of $r_t^2$ at lags 1 and 5
are too low in the model (median 0.184 and 0.133 against 0.242 and 0.191), and the simulated
moments vary much less across tickers than the data moments: the interquartile range of the
simulated lag-1 autocorrelation, [0.169, 0.198], lies inside that of the data, [0.131, 0.345].
The kurtosis is the clearest failure, with a median of 5.62 in the model against 13.61 in the
data. With $c\le\sqrt{1.5}$ the variance process is not volatile enough to generate the handful
of crash days that dominate the sample kurtosis of daily stock returns, and relaxing the bound
with a daily step would re-create days with zero variance; sub-daily QE steps would allow a larger
$c$ (\cref{sec:data-generators}).

\begin{table}[tbp]
  \centering
  \caption{Heston calibration by the simulated method of moments: data and simulated moments
    (E3, 100 tickers, \TRAIN{}).}
  \label{tab:e3-heston-moments}
  \small
  \fittable{e3_heston_moments}
  \tablenote{Data: moments of each ticker's \TRAIN{} returns (1258 days). Simulated: averages
    over the \TRAIN{}-length paths of the calibration at the fitted parameters (common random
    numbers; kurtosis: median over paths). Median and interquartile range across the 100 tickers; median of the absolute
    relative error and of the weighted deviation $\abs{t}$ (square root of the weight times the
    difference) per moment. The kurtosis (Pearson's, equal to 3 for a Gaussian law) is a
    diagnostic and is not among the targeted moments.}
\end{table}
